\chapter{经历组织、日记与认知图式}

\section{经历片段}

经历片段（Episode）组织连续经历流中的事件。它保存成员观察或认知的引用、顺序、角色、时间范围和边界解释。参与者变化、场景变化、任务状态变化、主题或因果断裂、明确完成事件和持续闲置都可以成为边界信号。

经历片段具有独立身份和修订。边界可以拆分、合并和重新分段；边界修订只改变事件组织，不改写原始观察事件或成员记忆。会话描述运行连续性，经历片段描述事件组织，二者不共享身份。

\section{日记}

日记（Journal）是对一段经历、阶段或时期的可修订叙事对象。日记引用经历片段、观察事件、记忆及其修订，并保存生产者和来源谱系。

日记用于保留较长时间尺度的经验上下文，支持回顾、自我叙事、检索和离线认知。日记不承担外部事实权威，也不是形成记忆的必经阶段。依赖来源发生实质修订、撤回、抑制或清除时，日记进入重新核验或形成新修订。

\section{认知图式}

认知图式（CognitiveSchema）表示主体从多个经历和认知中形成的可泛化结构。认知图式属于记忆域，并作为独立对象族维护稳定身份、修订、适用范围和生命周期。

一个认知图式至少表达：结构性主张、适用范围、支持认知、反例、边界案例和来源谱系。新的支持证据可以更新支持关系而不修改图式内容；只有结构解释、适用范围或边界发生变化时才创建内容修订。

认知图式支持拆分、合并、修订和撤回。反例不会自动删除图式；反例参与适用范围判断、冲突处理和后续修订。

\section{标签与关联证据}

标签（Tag）是稳定的语义线索或概念身份。标签身份与显示名称分离；相同字符串只能作为解析候选，不能自动证明两个标签具有同一语义身份。

关联证据（AssociationEvidence）记录认知对象之间有来源、有类型的关联依据，例如共现、顺序、显式关系、程序关系、共同结果或图式支持。关联证据保存来源、方向、极性、有效范围和生产者信息。

检索图中的中心性、传播质量和边权不构成认知对象。系统不定义同时承担认知结构和检索中心性的统一权威对象。认知结构由认知图式表达；检索拓扑由服务层根据当前有效关联证据重建。
